% Lösungen Prozentrechnung · Veränderung · Prüfungs-Fokus MSA (zusammenbau v0.6)
\einheitenkopf{Einheit 5 · Prozentuale Veränderung}

\erg{1}{a) auf $80\,\%$ gesenkt \quad b) $10\,\% = 8\,€$; $80 + 8 = 88\,€$ \quad c) $45 \cdot 1{,}2 = 54\,€$}
\erg{2}{a) $1{,}29 - 1{,}15 = 0{,}14$; $0{,}14 : 1{,}15 \approx 0{,}122 \approx 12{,}2\,\%$ \quad b) $1{,}90 - 1{,}70 = 0{,}20$; $0{,}20 : 1{,}70 \approx 0{,}118 \approx 11{,}8\,\%$ \quad c) $1\,340 - 1\,120 = 220$; $220 : 1\,340 \approx 0{,}164 \approx 16{,}4\,\%$ \quad d) $860 - 780 = 80$; $80 : 860 \approx 0{,}093 \approx 9{,}3\,\%$ \quad e) $57\,600 - 48\,000 = 9\,600$; $9\,600 : 48\,000 = 0{,}2 = 20\,\%$ (nicht durch den neuen Wert) \quad f) $72\,000 - 64\,000 = 8\,000$; $8\,000 : 64\,000 = 0{,}125 = 12{,}5\,\%$ (nicht durch den neuen Wert) \quad g) $2{,}40 \cdot 1{,}25 = 3{,}00\,€$ (nicht nur $0{,}60\,€$) \quad h) $65 \cdot 1{,}12 = 72{,}80\,€$ (nicht nur $7{,}80\,€$) \quad i) voller Preis $2 \cdot 14 + 2 \cdot 9 = 46\,€$; $46 \cdot 0{,}85 = 39{,}10\,€$ (nicht der Rabatt $6{,}90\,€$) \quad j) voller Preis $3 \cdot 22 + 13 = 79\,€$; $79 \cdot 0{,}75 = 59{,}25\,€$ (nicht der Rabatt $19{,}75\,€$) \quad k) Lücken: $100$ m und $5$ m; $15 : 240 = 0{,}0625 \approx 6{,}3\,\% > 5\,\%$ – Herr Lenz hat recht. \quad l) $15$ m; $1{,}1 : 8 = 0{,}1375 \approx 13{,}8\,\% < 15\,\%$ – Frau Roth hat nicht recht, die Zufahrt ist zulässig.}
